\documentclass[11pt]{article}
% SABLE-HE macro and package file.
% Keep this file small so the project compiles on a standard TeX Live setup.

\usepackage[T1]{fontenc}
\usepackage[utf8]{inputenc}
\usepackage{lmodern}
\usepackage{microtype}
\usepackage[a4paper,margin=1in]{geometry}
\usepackage{amsmath,amssymb,amsthm,mathtools}
\usepackage{booktabs,array,longtable}
\usepackage{enumitem}
\usepackage{xcolor}
\usepackage{url}
\usepackage[colorlinks=true,linkcolor=blue!50!black,citecolor=blue!50!black,urlcolor=blue!50!black]{hyperref}

\setlist[itemize]{leftmargin=1.2em}
\setlist[enumerate]{leftmargin=1.5em}

% Theorem environments.
\newtheorem{theorem}{Theorem}[section]
\newtheorem{lemma}[theorem]{Lemma}
\newtheorem{proposition}[theorem]{Proposition}
\newtheorem{corollary}[theorem]{Corollary}
\theoremstyle{definition}
\newtheorem{definition}[theorem]{Definition}
\newtheorem{remark}[theorem]{Remark}
\newtheorem{assumption}[theorem]{Assumption}
\newtheorem{construction}[theorem]{Construction}

% Algebra and crypto notation.
\newcommand{\F}{\mathbb{F}}
\newcommand{\Z}{\mathbb{Z}}
\newcommand{\N}{\mathbb{N}}
\newcommand{\poly}{\operatorname{poly}}
\newcommand{\negl}{\operatorname{negl}}
\newcommand{\supp}{\operatorname{supp}}
\newcommand{\wt}{\operatorname{wt}}
\newcommand{\Enc}{\mathsf{Enc}}
\newcommand{\Dec}{\mathsf{Dec}}
\newcommand{\KeyGen}{\mathsf{KeyGen}}
\newcommand{\Eval}{\mathsf{Eval}}
\newcommand{\Expand}{\mathsf{Expand}}
\newcommand{\Compact}{\mathsf{Compact}}
\newcommand{\EvalAdd}{\mathsf{EvalAdd}}
\newcommand{\EvalMul}{\mathsf{EvalMul}}
\newcommand{\sk}{\mathsf{sk}}
\newcommand{\ek}{\mathsf{ek}}
\newcommand{\ct}{\mathsf{ct}}
\newcommand{\Adv}{\mathsf{Adv}}
\newcommand{\LPN}{\mathsf{LPN}}
\newcommand{\SLPN}{\mathsf{sLPN}}
\newcommand{\CLPN}{\mathsf{CLPN}}
\newcommand{\GSW}{\mathsf{GSW}}
\newcommand{\Regev}{\mathsf{Regev}}
\newcommand{\Good}{\mathsf{Good}}
\newcommand{\Id}{\mathbf{I}}
\newcommand{\one}{\mathbf{1}}
\newcommand{\zero}{\mathbf{0}}
\newcommand{\calF}{\mathcal{F}}
\newcommand{\calA}{\mathcal{A}}
\newcommand{\calD}{\mathcal{D}}
\newcommand{\calC}{\mathcal{C}}
\newcommand{\getsr}{\stackrel{\$}{\leftarrow}}
\newcommand{\transpose}{^{\mathsf T}}
\newcommand{\inner}[2]{\langle #1,#2\rangle}
\newcommand{\til}[1]{\widetilde{#1}}

% Boxed algorithm helper.
\newcommand{\AlgTitle}[1]{\vspace{0.5em}\noindent\textbf{#1}\vspace{0.2em}}
\newenvironment{algobox}{\begin{center}\begin{minipage}{0.96\linewidth}\hrule\vspace{0.6em}}{\vspace{0.4em}\hrule\end{minipage}\end{center}}

% Project status labels.
\newcommand{\candidate}{\textsc{candidate}}
\newcommand{\sable}{\textsc{SABLE-HE}}


\title{\sable: A Sparse-LPN and Code-LPN Candidate for All-Code-Based Leveled Homomorphic Encryption}
\author{Research design draft\\Prepared as a mathematical construction and validation plan}
\date{Draft v0.9 C7 relation-resistant milestone -- \today}

\begin{document}
\maketitle

\begin{abstract}
This document proposes \sable, a candidate post-quantum leveled homomorphic encryption construction for bounded-depth arithmetic circuits over a prime field, and records the final v0.9-C7 relation-resistant validation pass.  The construction combines a sparse-LPN GSW-style nonlinear evaluation layer with an LPN/code-based linearly homomorphic compaction layer.  The purpose is assumption diversification: the design avoids Ring-LWE, Module-LWE, DDH, DCR, RSA, pairings, and Paillier-type assumptions, and instead relies on sparse-LPN and q-ary LPN/code-decoding assumptions.  We give formal syntax, algorithms, correctness invariants, a hybrid security proof under stated assumptions, efficiency formulas, parameter constraints, and a validation plan.  The result should be read as a research-grade \candidate{} construction, not as a deployed or peer-reviewed cryptographic standard.
\end{abstract}

\paragraph{Status and intended use.}
This is a theory-and-design manuscript.  It is suitable as the basis for a research paper, prototype library, security-estimation work, and cryptanalysis.  It does not certify concrete parameters as secure.  Any real deployment would require independent review, parameter estimation against current LPN and information-set-decoding attacks, side-channel review, and implementation audits.

\tableofcontents

\section{Introduction}

Homomorphic encryption (HE) allows a server to evaluate functions on encrypted data and return an encrypted result.  Modern practical HE is dominated by lattice assumptions, especially LWE, RLWE, and their ring/module variants.  This dominance is scientifically understandable: lattice HE has a mature lineage from Gentry's first FHE construction, BGV, GSW, BFV, CKKS, FHEW, and TFHE \cite{Gentry2009,BGV2014,GSW2013,FanVercauteren2012,CKKS2017,DucasMicciancio2015,CGGI2020}.  However, it also creates an assumption-diversity problem.  Post-quantum cryptography already benefits from several hardness families: lattices, codes, multivariate equations, hash-based signatures, and isogeny-style problems.  In contrast, post-quantum HE remains much more concentrated around lattices.

Recent work makes a code-based direction credible.  Corrigan-Gibbs, Henzinger, Kalai, and Vaikuntanathan construct somewhat homomorphic encryption from sparse LPN plus a linearly homomorphic encryption primitive \cite{CorriganGibbsHenzingerKalaiVaikuntanathan2024}.  Their work is the closest technical ancestor of this document.  It shows that sparse LPN can support bounded multiplication through a GSW-like approximate-eigenvector method.  Separately, Bitzer, Egger, Liu, and Wachter-Zeh propose post-quantum secure aggregation via code-based homomorphic encryption under LPN-style assumptions \cite{BitzerEggerLiuWachterZeh2026}.  A recent review also explicitly identifies code-based HE as an underexplored direction for post-quantum homomorphic encryption \cite{BhoiArakalaCormanRao2025}.  Random-code/rank-metric SHE has also recently been studied \cite{AguilarMelchorDyserynGaborit2025}.

The construction proposed here is called \sable, short for \emph{Sparse-LPN All-Code Bounded-depth Leveled Homomorphic Encryption}.  Its central idea is:
\[
\boxed{
\text{sparse-LPN GSW-style nonlinear evaluation}
\;+
\text{q-ary LPN/code-based linear compaction}
\;=
\text{all-code-based leveled SHE}.
}
\]
The sparse-LPN layer handles bounded multiplication.  The q-ary code/LPN layer homomorphically evaluates the final linear decryption map and compresses the output.  The proposed compaction layer is the main design change relative to sparse-LPN SHE constructions that instantiate linear homomorphism using non-code assumptions such as DDH or DCR.

\subsection{Research objective}

The objective is not to claim a complete full FHE scheme.  The objective is to design a mathematically coherent, non-lattice, post-quantum candidate for the following target functionality:
\[
\calF_{D,A,q} =
\left\{
 f:\F_q^m\to\F_q
 \;\middle|\;
 f \text{ is represented by an arithmetic circuit of multiplicative depth }D
 \text{ and addition budget }A
\right\}.
\]
Here $q\ge 3$ is prime.  Boolean inputs can be embedded as $0,1\in\F_q$.  Arithmetic multiplication gives AND, while OR, NOT, and XOR can be represented by low-degree polynomials:
\[
\operatorname{AND}(x,y)=xy,
\qquad
\operatorname{OR}(x,y)=x+y-xy,
\qquad
\operatorname{NOT}(x)=1-x,
\qquad
\operatorname{XOR}(x,y)=x+y-2xy.
\]
This makes \sable{} a candidate for private low-degree analytics, voting/counting with nonlinear checks, secure aggregation plus low-degree validation, private feature-interaction models, and small Boolean circuits.

\subsection{Contributions}

This draft makes four contributions.

\begin{enumerate}
    \item It formulates an all-code-based leveled SHE candidate using two independent LPN-style layers: sparse LPN for nonlinear evaluation and q-ary LPN/code encryption for linear compaction.
    \item It gives precise algorithms: $\KeyGen$, compact input $\Enc$, $\Expand$, $\EvalAdd$, $\EvalMul$, $\Compact$, and $\Dec$.
    \item It proves correctness by tracking a row-sparsity and coordinate-error invariant.  The invariant is inherited from the approximate-eigenvector technique of GSW and from sparse-LPN SHE analysis \cite{GSW2013,CorriganGibbsHenzingerKalaiVaikuntanathan2024}.
    \item It gives a hybrid IND-CPA security proof under explicit sparse-LPN and q-ary LPN assumptions, plus an attack and validation checklist for parameter selection.
\end{enumerate}

\subsection{Non-goals}

The following are intentionally outside the first version.
\begin{itemize}
    \item \textbf{No full bootstrapping claim.}  The construction is leveled and supports a predetermined bounded multiplicative depth.
    \item \textbf{No deployment-grade concrete parameters.}  Parameter tables in this document are formulas and illustrative examples only.  Security estimation must be done separately.
    \item \textbf{No CCA security.}  The security target is IND-CPA for a secret-key HE scheme with public evaluation keys.
    \item \textbf{No claim of superiority over lattice FHE.}  Mature lattice libraries remain the practical baseline.  The scientific contribution is assumption diversity and a code-based path to bounded homomorphism.
\end{itemize}

\subsection{Relationship to post-quantum standardization}

NIST's first post-quantum standards are for key establishment and signatures, not homomorphic encryption.  FIPS 203 specifies ML-KEM and relates its security to Module-LWE, while FIPS 204 specifies ML-DSA for digital signatures \cite{NISTFIPS203,NISTFIPS204}.  HE security guidance is maintained separately by the HomomorphicEncryption.org community \cite{AlbrechtEtAl2024HEGuidelines,HEStandard2018}.  Thus \sable{} should not be read as a NIST-standardized primitive.  It is a research proposal in the code-based HE family.

\section{Preliminaries}

\subsection{Notation}

Let $q\ge 3$ be prime and let $\F_q$ be the field of $q$ elements.  Vectors are column vectors unless stated otherwise.  For $x\in\F_q^n$, $\wt(x)$ denotes Hamming weight and $\supp(x)$ denotes support.  For $s\in\F_q^n$, define the extended secret
\[
    \tilde{s}=(-s_1,\ldots,-s_n,1)\in\F_q^{n+1}.
\]
For $c=(a,b)\in\F_q^n\times\F_q$, decryption-style inner products are written
\[
    c\cdot \tilde{s}=-a\transpose s+b.
\]
All probabilities are over the randomness of key generation, encryption, and noise unless otherwise specified.

\subsection{q-ary symmetric noise}

For $0\le \eta\le 1$, define the q-ary symmetric noise distribution $\chi_{q,\eta}$ over $\F_q$ by
\[
    e=0 \text{ with probability }1-\eta,
    \qquad
    e=u \text{ with probability }\eta/(q-1)\text{ for each }u\in\F_q^*.
\]
This is the natural q-ary analogue of Bernoulli noise for LPN.  It is also useful for coding-theoretic decoding because an error coordinate is either correct or replaced by a uniformly random nonzero field element.

\begin{lemma}[q-ary piling-up]
\label{lem:qary-piling}
Let $E_1,\ldots,E_B$ be independent q-ary symmetric errors with error probability $\eta$.  Let $\alpha_i\in\F_q^*$ be nonzero coefficients.  Then
\[
    E=\sum_{i=1}^B \alpha_iE_i
\]
is q-ary symmetric with error probability
\[
    \eta_B=\frac{q-1}{q}\left(1-\left(1-\frac{q\eta}{q-1}\right)^B\right).
\]
In particular, $\eta_B\le B\eta$ by the union bound.
\end{lemma}

\noindent The proof is included in Appendix~\ref{app:piling-up}.  The exact expression is useful for compaction parameters; the union bound is useful for simple correctness estimates.

\subsection{Sparse q-ary LPN}

Let $S_{n,k,q}$ be the distribution over $\F_q^n$ that samples a uniformly random support of size $k$ and assigns each supported position a uniformly random nonzero element of $\F_q^*$.  A sparse q-ary LPN sample with secret $s\in\F_q^n$ is
\[
    (a, a\transpose s+e),
    \qquad a\leftarrow S_{n,k,q},\quad e\leftarrow\chi_{q,\eta}.
\]

\begin{assumption}[Sparse q-ary LPN]
\label{ass:sparse-lpn}
For parameters $(n,k,q,\eta,M)$, the distribution of $M$ sparse q-ary LPN samples
\[
    \{(a_i,a_i\transpose s+e_i)\}_{i=1}^M
\]
with $s\leftarrow\F_q^n$, $a_i\leftarrow S_{n,k,q}$, and $e_i\leftarrow\chi_{q,\eta}$ is computationally indistinguishable from
\[
    \{(a_i,u_i)\}_{i=1}^M,
    \qquad u_i\leftarrow\F_q.
\]
\end{assumption}

Sparse LPN is a coding-theoretic noisy-decoding assumption.  Classical LPN and related noisy linear-equation assumptions are connected to learning parity with noise, random linear codes, and average-case coding hardness \cite{BlumKalaiWasserman2003,Alekhnovich2003,BerlekampMcEliecevanTilborg1978}.  Sparse-LPN variants have recently been used for homomorphic encryption \cite{CorriganGibbsHenzingerKalaiVaikuntanathan2024}.

\subsection{Dense q-ary LPN/code encryption}

The compaction layer uses a standard q-ary LPN/code primitive.  A dense q-ary LPN sample is
\[
    (A,Ar+e)
\]
for secret $r\in\F_q^{n_c}$, matrix $A\in\F_q^{m_c\times n_c}$, and error vector $e\in\F_q^{m_c}$ with independent $\chi_{q,\eta_c}$ coordinates.  Message encryption adds a codeword $\mathsf{Code}(x)$ to $Ar+e$.  In the simplest prototype, $\mathsf{Code}(x)=x\one$ is a repetition code.  A paper-grade version should replace repetition by a stronger q-ary error-correcting code.

\begin{assumption}[q-ary LPN/code pseudorandomness]
\label{ass:clpn}
For parameters $(n_c,m_c,q,\eta_c,M_c)$, polynomially many ciphertext blocks of the form
\[
    (A_j,A_jr+e_j+\mathsf{Code}(x_j))
\]
are semantically secure for messages $x_j\in\F_q$ under the q-ary LPN assumption and the chosen code/decoder.
\end{assumption}

This assumption is intentionally stated at the primitive level because the eventual library may instantiate $\mathsf{Code}$ by repetition, BCH-like q-ary codes, LDPC-like codes, or CRT-based code decompositions.

\subsection{Security target}

The target is secret-key IND-CPA security with public evaluation key.  The adversary may see the evaluation key, request polynomially many encryptions, choose two challenge messages, receive an encryption of one of them, and run arbitrary public homomorphic evaluation.  Since evaluation is deterministic and public, security follows from ciphertext and evaluation-key pseudorandomness under the stated assumptions.

\section{Design overview}

\sable{} uses three independent secrets:
\[
    t\in\F_q^n,\qquad s\in\F_q^n,
    \qquad r\in\F_q^{n_c}.
\]
The secret $t$ is used for compact input encryption.  The secret $s$ is used for GSW-style matrix ciphertexts and nonlinear evaluation.  The secret $r$ is used only for the code/LPN compaction layer.

The key chain is
\[
    t \longrightarrow s \longrightarrow r.
\]
An expansion key encrypts the coordinates of $\tilde{t}$ under the GSW-style secret $s$.  This turns compact sparse-LPN input ciphertexts into matrix ciphertexts.  A compaction key encrypts the coordinates of $\tilde{s}$ under the code/LPN compaction secret $r$.  This lets the evaluator homomorphically compute the final linear decryption form and return a compact output.

\subsection{Three ciphertext forms}

\paragraph{Input ciphertexts.}
A compact input ciphertext under $t$ has the form
\[
    c=(a,b)=(a,a\transpose t+\mu+e)\in\F_q^{n+1},
\]
where $a$ is $k$-sparse.  It satisfies
\[
    c\cdot\tilde{t}=\mu+e.
\]
This is small and suitable for data-owner encryption.

\paragraph{Evaluation ciphertexts.}
A GSW-style evaluation ciphertext under $s$ is a matrix $C\in\F_q^{N\times N}$, where $N=n+1$, satisfying
\[
    C\tilde{s}=\mu\tilde{s}+e.
\]
Homomorphic addition and multiplication are simply
\[
    C_1+C_2,
    \qquad
    C_1C_2.
\]
The price is that row support and error grow with circuit depth.

\paragraph{Compact output ciphertexts.}
After evaluation, if
\[
    C_*\tilde{s}=f(\mu_1,\ldots,\mu_m)\tilde{s}+e_*,
\]
then the last row $\rho_*$ satisfies
\[
    \rho_*\cdot\tilde{s}=f(\mu_1,\ldots,\mu_m)+e_{*,N}.
\]
The evaluator uses code/LPN encryptions of $\tilde{s}_i$ to homomorphically compute an encryption of $\rho_*\cdot\tilde{s}$.  This is the compaction step.

\subsection{Why the construction is code-based}

Both hardness assumptions are coding-theoretic.
\begin{itemize}
    \item Sparse LPN is a noisy sparse linear-code decoding problem.  It supplies compact input encryption and the GSW-style nonlinear layer.
    \item q-ary LPN/code encryption is a noisy random linear-code primitive.  It supplies linearly homomorphic compaction.
\end{itemize}
No lattice, number-theoretic, or pairing assumption appears in the construction.  This is the intended novelty.

\subsection{Main technical invariant}

The key invariant is that evaluation ciphertexts remain \emph{approximately eigenvector matrices}.  A ciphertext $C$ encrypting $\mu$ satisfies
\[
    C\tilde{s}=\mu\tilde{s}+e.
\]
Then
\[
    (C_1+C_2)\tilde{s}=(\mu_1+\mu_2)\tilde{s}+(e_1+e_2),
\]
and
\[
    C_1C_2\tilde{s}
    =C_1(\mu_2\tilde{s}+e_2)
    =\mu_1\mu_2\tilde{s}+\mu_2e_1+C_1e_2.
\]
Thus multiplication is valid if $C_1e_2$ remains sparse/noisy enough.  This is why row support matters.

\section{Core primitives}

Let $N=n+1$.  For $u\in\F_q^n$, write $\tilde{u}=(-u,1)\in\F_q^N$.

\subsection{Sparse-LPN zero encryption}

\begin{construction}[Sparse-LPN zero vector encryption]
\label{con:zero-vector}
For secret $u\in\F_q^n$, sample $a\leftarrow S_{n,k,q}$ and $e\leftarrow\chi_{q,\eta}$.  Output
\[
    z=(a,a\transpose u+e)\in\F_q^{N}.
\]
Then
\[
    z\cdot\tilde{u}=e.
\]
The row support is at most $k+1$.
\end{construction}

This is a noisy encryption of zero.  Adding a message in the last coordinate gives compact message encryption.

\subsection{Compact input encryption}

\begin{construction}[Compact sparse-LPN input encryption]
\label{con:input-enc}
For secret $t\in\F_q^n$ and message $\mu\in\F_q$:
\begin{enumerate}
    \item Sample $z=(a,a\transpose t+e)$ as in Construction~\ref{con:zero-vector}.
    \item Output
    \[
        c=z+\mu e_N=(a,a\transpose t+e+\mu),
    \]
    where $e_N=(0,\ldots,0,1)\in\F_q^N$.
\end{enumerate}
Then
\[
    c\cdot\tilde{t}=\mu+e.
\]
\end{construction}

The decryption error is not rounded away; it is corrected statistically by repetition and final plurality.  This is natural in the LPN/code setting.

\subsection{Sparse-LPN GSW-style zero matrix}

\begin{construction}[Sparse-LPN GSW-style zero matrix]
\label{con:gsw-zero}
For secret $s\in\F_q^n$, sample $N$ independent sparse-LPN zero rows
\[
    z_j=(a_j,a_j\transpose s+e_j),\qquad j=1,\ldots,N.
\]
Let $Z\in\F_q^{N\times N}$ be the matrix with rows $z_j$.  Then
\[
    Z\tilde{s}=e,
\]
where $e=(e_1,\ldots,e_N)\transpose$.  Every row has support at most $k+1$.
\end{construction}

\subsection{GSW-style scalar encryption}

\begin{construction}[Sparse-LPN GSW-style scalar encryption]
\label{con:gsw-scalar}
To encrypt $\alpha\in\F_q$ under $s$:
\begin{enumerate}
    \item Sample $Z$ as in Construction~\ref{con:gsw-zero}.
    \item Output
    \[
        X=Z+\alpha I_N.
    \]
\end{enumerate}
Then
\[
    X\tilde{s}=\alpha\tilde{s}+e.
\]
Each row has support at most $k+2$.
\end{construction}

This formulation is important for security: $X$ is a pseudorandom zero-encryption matrix plus a public diagonal shift.  This mirrors the approximate-eigenvector technique in GSW-style encryption \cite{GSW2013} and its sparse-LPN adaptation \cite{CorriganGibbsHenzingerKalaiVaikuntanathan2024}.

\subsection{Expansion from compact input to GSW form}

The expansion key contains GSW encryptions of the coordinates of $\tilde{t}$:
\[
    \mathsf{EK}^{\rm exp}_i=\GSW\Enc_s(\tilde{t}_i),
    \qquad i=1,\ldots,N.
\]
Given compact ciphertext $c=(c_1,\ldots,c_N)$, define
\[
    \Expand(c)=\sum_{i=1}^{N}c_i\mathsf{EK}^{\rm exp}_i.
\]
Since $c$ is sparse with at most $k+1$ nonzero coordinates, expansion sums at most $k+1$ GSW matrices.

\begin{lemma}[Expansion identity]
\label{lem:expansion-identity}
Let $c$ be a compact encryption of $\mu$ under $t$.  Let $C=\Expand(c)$.  Then
\[
    C\tilde{s}=\mu\tilde{s}+e',
\]
where
\[
    e'=e_0\tilde{s}+\sum_{i\in\supp(c)}c_i e^{(i)}.
\]
Here $e_0=c\cdot\tilde{t}-\mu$ is the input error and $e^{(i)}$ is the GSW error vector in $\mathsf{EK}^{\rm exp}_i$.
\end{lemma}

\begin{proof}
For each $i$,
\[
    \mathsf{EK}^{\rm exp}_i\tilde{s}=\tilde{t}_i\tilde{s}+e^{(i)}.
\]
Therefore
\[
    C\tilde{s}
    =\sum_i c_i\tilde{t}_i\tilde{s}+\sum_i c_ie^{(i)}
    =(c\cdot\tilde{t})\tilde{s}+\sum_i c_ie^{(i)}.
\]
Since $c\cdot\tilde{t}=\mu+e_0$, the claim follows.
\end{proof}

\subsection{Code-LPN linearly homomorphic encryption}

Let $\mathsf{Code}:\F_q\to\F_q^{m_c}$ be a linear error-correcting encoding.  For a first prototype, $\mathsf{Code}(x)=x\one$.  Let $\mathsf{Decode}$ be its decoder.

\begin{construction}[CLPN encryption]
\label{con:clpn}
For secret $r\in\F_q^{n_c}$ and message $x\in\F_q$:
\begin{enumerate}
    \item Sample $A\leftarrow\F_q^{m_c\times n_c}$.
    \item Sample $e\in\F_q^{m_c}$ with independent $\chi_{q,\eta_c}$ coordinates.
    \item Output
    \[
        Z=(A,b=Ar+\mathsf{Code}(x)+e).
    \]
\end{enumerate}
Decryption computes
\[
    y=b-Ar=\mathsf{Code}(x)+e
\]
and outputs $\mathsf{Decode}(y)$.
\end{construction}

\begin{construction}[CLPN linear evaluation]
Given ciphertexts $Z_i=(A_i,b_i)$ encrypting $x_i$ and coefficients $\alpha_i\in\F_q$, output
\[
    \sum_i \alpha_i Z_i
    =\left(\sum_i\alpha_iA_i,\sum_i\alpha_i b_i\right).
\]
This decrypts to
\[
    \mathsf{Code}\left(\sum_i\alpha_ix_i\right)+\sum_i\alpha_ie_i.
\]
\end{construction}

\section{Full construction}

This section defines the \sable{} scheme.  The scheme is secret-key encryption with a public evaluation key.

\subsection{Parameters}

A parameter set is
\[
    \Pi=(q,n,k,\eta,D,A,R,n_c,m_c,\eta_c,\mathsf{Code}).
\]
The parameters mean:
\begin{center}
\begin{tabular}{ll}
\toprule
Parameter & Meaning\\
\midrule
$q$ & prime plaintext/noise field modulus, $q\ge 3$\\
$n$ & sparse-LPN secret dimension\\
$k$ & sparse row weight\\
$\eta$ & sparse-LPN q-ary error probability\\
$D$ & supported multiplicative depth\\
$A$ & conservative addition/monomial budget\\
$R$ & number of independent replicas for failure amplification\\
$n_c,m_c$ & compactor secret length and code block length\\
$\eta_c$ & compactor q-ary error probability\\
$\mathsf{Code}$ & linear q-ary error-correcting code for compaction\\
\bottomrule
\end{tabular}
\end{center}
Let $N=n+1$.

\subsection{KeyGen}

\begin{algobox}
\AlgTitle{$\KeyGen$}
\begin{enumerate}
    \item Choose parameters $\Pi$ for security parameter $\lambda$ and depth $D$.
    \item Sample independent secrets
    \[
        t\leftarrow\F_q^n,\qquad s\leftarrow\F_q^n,
        \qquad r\leftarrow\F_q^{n_c}.
    \]
    \item Define $\tilde{t}=(-t,1)$ and $\tilde{s}=(-s,1)$.
    \item For each $i\in[N]$, generate an expansion-key matrix
    \[
        \mathsf{EK}^{\rm exp}_i=\GSW\Enc_s(\tilde{t}_i).
    \]
    \item For each $i\in[N]$, generate a compaction-key ciphertext
    \[
        \mathsf{EK}^{\rm cmp}_i=\CLPN\Enc_r(\tilde{s}_i).
    \]
    \item Output
    \[
        \sk=(t,s,r),
        \qquad
        \ek=(\Pi,\mathsf{EK}^{\rm exp},\mathsf{EK}^{\rm cmp}).
    \]
\end{enumerate}
\end{algobox}

\subsection{Enc}

For a message $\mu\in\F_q$, generate $R$ independent compact ciphertexts:
\[
    c^{(j)}=\Regev\Enc_t(\mu),\qquad j=1,\ldots,R.
\]
Return
\[
    \ct=(c^{(1)},\ldots,c^{(R)}).
\]
Replicas are used only to reduce the decryption-failure probability.  They may be omitted in toy experiments.

\subsection{Expand}

For each replica,
\[
    C^{(j)}=\sum_{i=1}^{N}c_i^{(j)}\mathsf{EK}^{\rm exp}_i.
\]
Return
\[
    \mathsf{CT}=(C^{(1)},\ldots,C^{(R)}).
\]
The result is a vector of GSW-style matrix ciphertexts under secret $s$.

\subsection{Homomorphic evaluation}

For expanded ciphertexts
\[
    \mathsf{CT}_1=(C_1^{(1)},\ldots,C_1^{(R)}),
    \qquad
    \mathsf{CT}_2=(C_2^{(1)},\ldots,C_2^{(R)}),
\]
define
\[
    \EvalAdd(\mathsf{CT}_1,\mathsf{CT}_2)=
    \left(C_1^{(j)}+C_2^{(j)}\right)_{j=1}^{R},
\]
and
\[
    \EvalMul(\mathsf{CT}_1,\mathsf{CT}_2)=
    \left(C_1^{(j)}C_2^{(j)}\right)_{j=1}^{R}.
\]
Multiplication may be oriented so that the matrix with smaller row support is placed on the left, because the error bound for multiplication is asymmetric.

\subsection{Compact}

For each evaluated replica $C_*^{(j)}$:
\begin{enumerate}
    \item Let $\rho_*^{(j)}\in\F_q^N$ be the last row of $C_*^{(j)}$.
    \item Homomorphically evaluate the linear form $x\mapsto\rho_*^{(j)}\cdot x$ on the compaction-key encryptions of $\tilde{s}_i$:
    \[
        Z^{(j)}=
        \mathsf{CLPN.EvalLin}
        \left(\rho_*^{(j)};
        \mathsf{EK}^{\rm cmp}_1,\ldots,\mathsf{EK}^{\rm cmp}_N\right).
    \]
\end{enumerate}
Return
\[
    \mathsf{ct}_{\rm out}=(Z^{(1)},\ldots,Z^{(R)}).
\]

\subsection{Dec}

For each replica, compute
\[
    y^{(j)}=\CLPN\Dec_r(Z^{(j)}).
\]
Return the plurality value among $y^{(1)},\ldots,y^{(R)}$.

\subsection{Correctness intuition}

If evaluation gives
\[
    C_*^{(j)}\tilde{s}=f(\mu_1,\ldots,\mu_m)\tilde{s}+e_*^{(j)},
\]
then the last row $\rho_*^{(j)}$ satisfies
\[
    \rho_*^{(j)}\cdot\tilde{s}=f(\mu_1,\ldots,\mu_m)+e_{*,N}^{(j)}.
\]
The compaction key encrypts the coordinates of $\tilde{s}$.  Therefore $\Compact$ produces a code/LPN encryption of the decryption value.  If both the sparse-LPN evaluation error and the code/LPN compaction error are decoded correctly, decryption returns $f(\mu_1,\ldots,\mu_m)$.

\section{Correctness analysis}

\subsection{Good evaluation ciphertexts}

\begin{definition}[Good GSW-style ciphertext]
Let $C\in\F_q^{N\times N}$.  We say that $C$ is $(w,\varepsilon)$-good for $\mu\in\F_q$ under secret $s$ if:
\begin{enumerate}
    \item every row of $C$ has Hamming support at most $w$;
    \item writing
    \[
        C\tilde{s}=\mu\tilde{s}+e,
    \]
    every coordinate satisfies
    \[
        \Pr[e_j\ne 0]\le \varepsilon.
    \]
\end{enumerate}
\end{definition}

This definition tracks only row support and coordinatewise error probability.  It deliberately does not require independence between all error coordinates after evaluation.  The proofs use union bounds, so this weaker invariant is enough.

\subsection{Fresh expansion quality}

Let
\[
    w_0=(k+1)(k+2),
    \qquad
    \varepsilon_0=(k+2)\eta.
\]

\begin{lemma}[Fresh expanded ciphertext quality]
\label{lem:fresh-quality}
For a compact encryption $c=\Regev\Enc_t(\mu)$, the expanded ciphertext $C=\Expand(c)$ is $(w_0,\varepsilon_0)$-good for $\mu$.
\end{lemma}

\begin{proof}
The compact ciphertext $c$ has at most $k+1$ nonzero coordinates: at most $k$ from $a$ and one final coordinate.  Each expansion-key matrix has row support at most $k+2$ because it is a sparse-LPN zero matrix plus a diagonal shift.  Therefore each row of $C=\sum_i c_i\mathsf{EK}^{\rm exp}_i$ has support at most $(k+1)(k+2)=w_0$.

By Lemma~\ref{lem:expansion-identity},
\[
    C\tilde{s}=\mu\tilde{s}+e_0\tilde{s}+\sum_{i\in\supp(c)}c_ie^{(i)}.
\]
For a fixed coordinate $j$, the term $e_0\tilde{s}_j$ is nonzero only if $e_0\ne 0$ and $\tilde{s}_j\ne0$, so it is nonzero with probability at most $\eta$.  The summation contains at most $k+1$ independent GSW error coordinates, each nonzero with probability $\eta$.  A union bound gives
\[
    \Pr[e'_j\ne0]\le \eta+(k+1)\eta=(k+2)\eta.
\]
\end{proof}

\subsection{Addition and multiplication invariants}

\begin{lemma}[Homomorphic addition]
\label{lem:add-good}
If $C_1$ is $(w_1,\varepsilon_1)$-good for $\mu_1$ and $C_2$ is $(w_2,\varepsilon_2)$-good for $\mu_2$, then $C_1+C_2$ is $(w_1+w_2,\varepsilon_1+\varepsilon_2)$-good for $\mu_1+\mu_2$.
\end{lemma}

\begin{proof}
Row support is subadditive.  Also,
\[
    (C_1+C_2)\tilde{s}
    =C_1\tilde{s}+C_2\tilde{s}
    =(\mu_1+\mu_2)\tilde{s}+e_1+e_2.
\]
For each coordinate, $e_{1,j}+e_{2,j}$ is nonzero only if at least one of $e_{1,j},e_{2,j}$ is nonzero, except for possible cancellation.  Thus
\[
    \Pr[e_{1,j}+e_{2,j}\ne0]
    \le \varepsilon_1+\varepsilon_2.
\]
\end{proof}

\begin{lemma}[Homomorphic multiplication]
\label{lem:mul-good}
If $C_1$ is $(w_1,\varepsilon_1)$-good for $\mu_1$ and $C_2$ is $(w_2,\varepsilon_2)$-good for $\mu_2$, then $C_1C_2$ is $(w_1w_2,\varepsilon_1+w_1\varepsilon_2)$-good for $\mu_1\mu_2$.
\end{lemma}

\begin{proof}
A row of $C_1C_2$ is a linear combination of at most $w_1$ rows of $C_2$, each of support at most $w_2$.  Hence the row support is at most $w_1w_2$.

For correctness,
\[
    C_2\tilde{s}=\mu_2\tilde{s}+e_2,
\]
so
\[
    C_1C_2\tilde{s}=C_1(\mu_2\tilde{s}+e_2)=\mu_2C_1\tilde{s}+C_1e_2.
\]
Using $C_1\tilde{s}=\mu_1\tilde{s}+e_1$, we get
\[
    C_1C_2\tilde{s}=\mu_1\mu_2\tilde{s}+\mu_2e_1+C_1e_2.
\]
For a fixed output coordinate $j$, the term $(\mu_2e_1)_j$ is nonzero with probability at most $\varepsilon_1$.  The $j$th coordinate of $C_1e_2$ depends on at most $w_1$ coordinates of $e_2$, each nonzero with probability at most $\varepsilon_2$.  By a union bound,
\[
    \Pr[(C_1e_2)_j\ne0]\le w_1\varepsilon_2.
\]
Combining the two bad events gives the claimed bound.
\end{proof}

\subsection{Depth bound}

For a balanced multiplication tree of depth $D$ starting from fresh expanded ciphertexts, define recursively
\[
    w_{r+1}=w_r^2,
    \qquad
    \varepsilon_{r+1}\le (1+w_r)\varepsilon_r,
\]
with $w_0=(k+1)(k+2)$ and $\varepsilon_0=(k+2)\eta$.  Hence
\[
    w_D=w_0^{2^D}.
\]
A conservative closed-form bound is
\[
    \varepsilon_D
    \le
    \varepsilon_0\prod_{r=0}^{D-1}(1+w_r)
    \le
    (k+2)\eta\prod_{r=0}^{D-1}(1+w_0^{2^r}).
\]
For $w_0\ge2$, this is roughly
\[
    \varepsilon_D=O\left(\eta\,w_0^{2^D-1}\right).
\]
If at most $A$ additive branches are combined, a conservative final bound is
\[
    \varepsilon_f\le A(k+2)\eta\prod_{r=0}^{D-1}(1+w_0^{2^r}).
\]

\subsection{Compaction correctness}

Let $B$ be the number of nonzero entries in the final row $\rho_*$ used for compaction.  For a repetition-code compactor, $B$ independent nonzero scalar multiples of q-ary symmetric errors yield effective error probability
\[
    \eta_{c,B}=\frac{q-1}{q}\left(1-\left(1-\frac{q\eta_c}{q-1}\right)^B\right).
\]
For a general code $\mathsf{Code}$, let
\[
    \varepsilon_{\rm cmp}(B)
\]
be the decoder's failure probability at this effective error rate and block length $m_c$.

\begin{theorem}[One-replica correctness]
\label{thm:one-replica-correctness}
Suppose $C_*$ is $(w,\varepsilon)$-good for $\mu_*=f(\mu_1,\ldots,\mu_m)$ and the final row has support $B\le w$.  Then one compacted replica decrypts correctly with probability at least
\[
    1-\varepsilon-\varepsilon_{\rm cmp}(B).
\]
\end{theorem}

\begin{proof}
Write
\[
    C_*\tilde{s}=\mu_*\tilde{s}+e_*.
\]
Let $\rho_*$ be the last row.  Since $\tilde{s}_N=1$,
\[
    \rho_*\cdot\tilde{s}=\mu_*+e_{*,N}.
\]
By goodness, $\Pr[e_{*,N}\ne0]\le\varepsilon$.  If $e_{*,N}=0$, the compactor homomorphically evaluates the linear form $\rho_*\cdot x$ on encryptions of $\tilde{s}_i$ and therefore decrypts to $\mu_*$ unless the code decoder fails.  The decoder failure probability is $\varepsilon_{\rm cmp}(B)$.  A union bound proves the theorem.
\end{proof}

\begin{theorem}[Replicated correctness]
\label{thm:replicated-correctness}
Let
\[
    p=\varepsilon+\varepsilon_{\rm cmp}(B)<1/2.
\]
With $R$ independent replicas and plurality decoding, the final failure probability is at most
\[
    \exp\left(-2R\left(\frac12-p\right)^2\right).
\]
\end{theorem}

\begin{proof}
Treat each replica as a Bernoulli trial that fails with probability at most $p$.  Majority/plurality fails only if at least half the replicas fail.  Hoeffding's inequality gives the stated bound.  For $q>2$, plurality failure is at most the majority-failure event in which the correct value appears in fewer than half the replicas.
\end{proof}

\section{Security analysis}

The security claim is conditional and modular.  The construction is secure if the sparse-LPN and q-ary LPN/code distributions used in its layers are pseudorandom for the chosen parameter sizes and number of samples.  Concrete parameters require separate estimation.

\subsection{Pseudorandomness of zero encryptions}

The following proposition is immediate from Assumption~\ref{ass:sparse-lpn}.

\begin{proposition}[Sparse-LPN zero ciphertext pseudorandomness]
Sparse-LPN zero vectors of the form
\[
    (a,a\transpose u+e)
\]
are indistinguishable from sparse-supported vectors whose final coordinate is uniform in $\F_q$.  Likewise, sparse-LPN zero matrices $Z$ sampled rowwise are indistinguishable from matrices with the same row-support distribution and uniform final row entries.
\end{proposition}

Because $\GSW\Enc_s(\alpha)=Z+\alpha I_N$, semantic security of GSW-style scalar encryption follows by adding a fixed public shift to a pseudorandom matrix.  If the encrypted scalar is secret but independent of the GSW secret, the same argument applies after conditioning on that scalar.

\subsection{IND-CPA theorem}

\begin{theorem}[IND-CPA security of \sable]
\label{thm:security}
Assume:
\begin{enumerate}
    \item sparse q-ary LPN pseudorandomness for all compact input ciphertext samples under $t$;
    \item sparse q-ary LPN pseudorandomness for all GSW-style zero rows under $s$ in the expansion key;
    \item q-ary LPN/code semantic security for all compaction-key ciphertexts under $r$.
\end{enumerate}
Then \sable{} is IND-CPA secure for polynomially many encryption queries and polynomial-size public evaluation keys.
\end{theorem}

\begin{proof}
Let $\calA$ be a polynomial-time IND-CPA adversary.  We define hybrids.

\paragraph{Hybrid H0: real experiment.}
The adversary receives the real evaluation key
\[
    \ek=(\mathsf{EK}^{\rm exp},\mathsf{EK}^{\rm cmp})
\]
and oracle encryptions under $t$.  The challenge ciphertext encrypts either $\mu_0$ or $\mu_1$ under $t$.

\paragraph{Hybrid H1: replace the compaction key.}
Replace each compaction-key ciphertext
\[
    \mathsf{EK}^{\rm cmp}_i=\CLPN\Enc_r(\tilde{s}_i)
\]
by an encryption of zero, or equivalently by a pseudorandom q-ary LPN/code ciphertext.  The messages $\tilde{s}_i$ may depend on $s$ but are independent of the compaction secret $r$.  By Assumption~\ref{ass:clpn}, the adversary's distinguishing advantage changes by at most a negligible quantity.

\paragraph{Hybrid H2: replace the expansion key.}
Replace each expansion-key matrix
\[
    \mathsf{EK}^{\rm exp}_i=\GSW\Enc_s(\tilde{t}_i)=Z_i+\tilde{t}_iI_N
\]
by a pseudorandom matrix with the same public support profile.  Conditioned on $t$, each $\tilde{t}_iI_N$ is a fixed shift and $Z_i$ is a sparse-LPN zero matrix under secret $s$.  By Assumption~\ref{ass:sparse-lpn}, replacing all such matrices changes the adversary's advantage negligibly.

\paragraph{Hybrid H3: replace encryption-oracle outputs.}
Now the evaluation key reveals no computational information about $t$.  Replace all compact input ciphertexts, including the challenge ciphertext, by sparse-supported pseudorandom vectors.  Each compact encryption has the form
\[
    (a,a\transpose t+e+\mu).
\]
For any fixed $\mu$, this is a fixed shift of a sparse-LPN sample.  By Assumption~\ref{ass:sparse-lpn}, encryptions of $\mu_0$ and $\mu_1$ are indistinguishable from pseudorandom vectors and hence from each other.

In Hybrid H3 the adversary's view is independent of the challenge bit.  Therefore its advantage is zero in H3, and its advantage in H0 is bounded by the sum of the negligible hybrid gaps.
\end{proof}

\subsection{What the proof does and does not cover}

The proof covers passive IND-CPA privacy for the abstract construction.  It does not cover:
\begin{itemize}
    \item chosen-ciphertext attacks;
    \item maliciously generated public parameters;
    \item side channels or timing leakage;
    \item concrete attack costs for a particular parameter set;
    \item algebraic attacks exploiting unexpected structure in an optimized implementation;
    \item correctness failures that leak through repeated decryption attempts.
\end{itemize}
These issues must be addressed before any implementation is considered deployment-grade.

\subsection{Relevant attack families}

The most relevant attacks for concrete parameter selection include:
\begin{itemize}
    \item BKW-style LPN attacks \cite{BlumKalaiWasserman2003};
    \item information-set decoding and Prange/Stern-style decoding attacks \cite{Prange1962,Stern1989};
    \item low-noise LPN attacks when $\eta$ or $\eta_c$ is extremely small;
    \item sample-amplification attacks, because the evaluation key contains many LPN samples;
    \item distinguishing attacks against sparse public rows;
    \item attacks exploiting the chain of encrypted secrets $t\to s\to r$.
\end{itemize}
The construction deliberately avoids circular encryption of a secret under itself.  The expansion key encrypts $\tilde{t}$ under independent secret $s$, and the compaction key encrypts $\tilde{s}$ under independent secret $r$.

\section{Efficiency and parameter constraints}

This section gives symbolic efficiency formulas.  It does not certify concrete security.  A real submission should include a parameter estimator and independent cryptanalysis.

\subsection{Ciphertext and key sizes}

Let $N=n+1$ and $w_0=(k+1)(k+2)$.

\begin{center}
\begin{tabular}{lll}
\toprule
Object & Dense size & Sparse/structured size\\
\midrule
Compact input ciphertext & $N$ field elements & $O(k\log n+\log q)$ bits plus values\\
One GSW matrix ciphertext & $N^2$ field elements & $O(Nw)$ field/index pairs\\
Expansion key & $N$ GSW matrices & $O(N^2(k+2))$ field/index pairs\\
Compaction key & $N$ CLPN ciphertexts & $O(Nm_cn_c)$ dense, seedable in practice\\
Compacted output & $R$ CLPN ciphertexts & independent of circuit size\\
\bottomrule
\end{tabular}
\end{center}

The expansion key is large but amortized.  It is published once per key and reused for many input ciphertexts.  The final compact output does not grow with the evaluated circuit.

\subsection{Evaluation costs}

For row-sparse matrices with row supports $w_1,w_2$:
\[
    \EvalAdd: O(N(w_1+w_2)),
    \qquad
    \EvalMul: O(Nw_1w_2)
\]
field operations, ignoring sorting/merging overhead.  After depth $D$,
\[
    w_D=w_0^{2^D}.
\]
Thus \sable{} is a low-depth scheme.  It is realistic to target degree-2, degree-4, or possibly degree-8 computations before considering further compression or bootstrapping.

\subsection{Correctness budget}

A conservative correctness target is
\[
    \varepsilon_f+\varepsilon_{\rm cmp}(B)<\frac12.
\]
A stronger practical target is
\[
    \varepsilon_f+\varepsilon_{\rm cmp}(B)\le 0.1,
\]
so that a moderate replica count $R$ gives negligible final failure.  Using the bound from Section~6,
\[
    \varepsilon_f
    \le
    A(k+2)\eta\prod_{r=0}^{D-1}(1+w_0^{2^r}).
\]
Therefore a simple design inequality is
\[
    \eta
    \le
    \frac{\tau}{A(k+2)\prod_{r=0}^{D-1}(1+w_0^{2^r})}
\]
for target sparse-LPN evaluation error $\tau$.

For the compactor, if the final row support is $B$, the effective q-ary error rate is
\[
    \eta_{c,B}=\frac{q-1}{q}\left(1-\left(1-\frac{q\eta_c}{q-1}\right)^B\right).
\]
The code should be chosen so that decoding at this error rate fails with probability at most the desired $\varepsilon_{\rm cmp}(B)$.

\subsection{Illustrative non-secure toy parameters}

The following parameters are only for debugging a Python prototype:
\[
    q=127,\quad n=64,\quad k=3,\quad D=1,\quad R=9.
\]
They are intentionally small and not secure.  Their purpose is to validate algebraic correctness and serialization.

A research-grade estimator must choose $n,k,q,\eta,n_c,m_c,\eta_c$ jointly.  The estimator should output:
\begin{itemize}
    \item estimated classical and quantum cost of sparse-LPN distinguishing;
    \item estimated cost of dense q-ary LPN/code distinguishing;
    \item estimated ISD cost for the relevant code parameters;
    \item total number of LPN samples exposed by input ciphertexts and evaluation keys;
    \item correctness failure estimates before and after replication.
\end{itemize}

\subsection{Parameter-selection recipe}

A practical research workflow is:
\begin{enumerate}
    \item Fix the application class, especially multiplicative depth $D$ and addition budget $A$.
    \item Choose a small sparse row weight $k$ that keeps multiplication feasible.
    \item Compute $w_0=(k+1)(k+2)$ and the depth growth $w_D=w_0^{2^D}$.
    \item Choose $\eta$ small enough for correctness but not so small that low-noise LPN attacks become cheap.
    \item Choose $q$ to balance message space, q-ary error behavior, and implementation simplicity.
    \item Choose $\mathsf{Code}$ and $(n_c,m_c,\eta_c)$ so that compaction decoding succeeds at row support $B\le\min(N,w_D)$.
    \item Run attack estimators and increase dimensions until the desired security category is reached.
\end{enumerate}

\subsection{Benchmark baselines}

The natural baselines are TFHE/FHEW for Boolean circuits, BFV/BGV for exact arithmetic, and CKKS for approximate arithmetic \cite{DucasMicciancio2015,CGGI2020,FanVercauteren2012,BGV2014,CKKS2017}.  \sable{} should not be expected to dominate those mature lattice schemes broadly.  It should be evaluated on the specific workloads where code-based assumption diversity and low-degree structure matter.

\section{Validation plan}

\sable{} should be validated in four parallel tracks: formal proofs, parameter estimation, prototype implementation, and benchmark comparison.

\subsection{Formal proof track}

The formal proof track should produce a machine-checkable or at least reviewer-checkable proof package for:
\begin{enumerate}
    \item sparse-LPN zero-vector pseudorandomness;
    \item GSW-style scalar encryption security as zero encryption plus diagonal shift;
    \item expansion correctness;
    \item row-support and coordinate-error invariants under addition and multiplication;
    \item code/LPN compaction correctness;
    \item hybrid IND-CPA security.
\end{enumerate}
The proof should explicitly track sample counts.  This matters because the evaluation key exposes $O(N^2)$ sparse-LPN rows and the compaction key exposes $O(Nm_c)$ dense q-ary LPN samples.

\subsection{Parameter-estimation track}

The estimator should be a separate, versioned artifact.  It should compute at least:
\begin{itemize}
    \item sparse-LPN distinguishing cost for $(n,k,q,\eta,M)$;
    \item q-ary LPN/code distinguishing cost for $(n_c,m_c,q,\eta_c,M_c)$;
    \item information-set decoding estimates for the compaction code;
    \item low-noise attack costs as $\eta$ and $\eta_c$ decrease;
    \item correctness failure from the exact q-ary piling-up formula;
    \item final failure after $R$ replicas.
\end{itemize}
No parameter set should be described as secure until this estimator and independent review agree.

\subsection{Prototype track}

The first Python library should implement the following objects:
\begin{itemize}
    \item finite-field arithmetic modulo a prime $q$;
    \item sparse vectors and row-sparse matrices;
    \item q-ary symmetric noise sampling;
    \item compact sparse-LPN input encryption;
    \item GSW-style sparse-LPN matrices;
    \item expansion, addition, multiplication, and compaction;
    \item repetition-code decoding for initial tests;
    \item a pluggable code interface for BCH/LDPC/CRT-style compaction later.
\end{itemize}

Core unit tests should verify:
\begin{enumerate}
    \item $c\cdot\tilde{t}=\mu+e$ for compact input ciphertexts;
    \item $X\tilde{s}=\alpha\tilde{s}+e$ for GSW-style encryptions;
    \item $\Expand(c)\tilde{s}=\mu\tilde{s}+e'$;
    \item addition and multiplication decrypt to correct low-degree arithmetic values with expected failure rates;
    \item compaction returns the same value as direct last-row decryption;
    \item empirical failure rates match the formulas in Section~6.
\end{enumerate}

\subsection{Benchmark track}

Benchmark on workloads that match the design:
\begin{itemize}
    \item degree-2 private feature interaction:
    \[
        f(x)=\sum_i a_ix_i+\sum_{i<j}b_{ij}x_ix_j;
    \]
    \item encrypted voting with polynomial validity checks;
    \item private counting with pairwise fraud/anomaly indicators;
    \item small Boolean circuits represented arithmetically over $\F_q$;
    \item secure aggregation plus low-degree validation.
\end{itemize}
Metrics should include input ciphertext size, expansion-key size, compaction-key size, expansion time, multiplication time, compaction time, final ciphertext size, empirical failure probability, and estimated security bits.

\subsection{Success criteria for a publishable paper}

A strong paper based on \sable{} should provide:
\begin{enumerate}
    \item a clean theorem statement under standard or carefully justified assumptions;
    \item an implementation that is small enough for independent reproduction;
    \item security estimates against known LPN and decoding attacks;
    \item a comparison against sparse-LPN SHE with number-theoretic LHE compaction;
    \item a comparison against TFHE/FHEW on small Boolean workloads;
    \item evidence that the all-code compactor is not merely aesthetic but changes assumption structure meaningfully.
\end{enumerate}

\section{Limitations and future work}

\subsection{Limitations}

The construction has several serious limitations.

\paragraph{Depth growth.}
Row support grows as $w_D=w_0^{2^D}$.  This restricts \sable{} to low-depth computation unless a new compression or bootstrapping method is added.

\paragraph{Low-noise tension.}
Correctness prefers small $\eta$ and $\eta_c$, but security against low-noise LPN attacks may prefer larger noise or larger dimensions.  This tension is central and must be quantified.

\paragraph{Large evaluation keys.}
The expansion key contains $N$ GSW matrices, each with $N$ sparse rows.  The compaction key also contains $N$ code/LPN ciphertexts.  Seeded randomness and structured sampling may reduce storage, but those optimizations need separate security analysis.

\paragraph{Secret-key setting.}
This draft gives a secret-key HE scheme with public evaluation keys.  Public-key encryption is not included.  A public-key variant may be possible using standard LPN public-key techniques, but it would need a new proof.

\paragraph{Compactor quality.}
A repetition-code compactor is simple but inefficient.  A serious version should use stronger q-ary codes and possibly CRT decomposition, inspired by recent code-based secure aggregation work \cite{BitzerEggerLiuWachterZeh2026}.

\subsection{Future work}

The most important next steps are:
\begin{enumerate}
    \item design a stronger q-ary code compactor with efficient decoding and exact failure estimates;
    \item build a sparse-LPN parameter estimator specialized to the exposed sample structure;
    \item add seeded matrix generation to shrink evaluation keys;
    \item study whether multiplication can be followed by a code-based refresh operation;
    \item explore batching or packing across rows without introducing ring/lattice assumptions;
    \item construct a public-key variant;
    \item implement a Python prototype and then a C/Rust reference implementation.
\end{enumerate}

\subsection{Recommended claim language}

For a research paper, use cautious claim language:

\begin{quote}
We propose a candidate all-code-based leveled somewhat homomorphic encryption scheme for bounded-depth arithmetic circuits.  Under sparse-LPN and q-ary LPN/code assumptions, the construction achieves compact input encryption, GSW-style nonlinear evaluation, and code/LPN-based output compaction without lattice or number-theoretic assumptions.
\end{quote}

Avoid stronger claims such as ``practical full FHE from codes'' or ``secure against all quantum attacks.''  The correct claim is conditional post-quantum plausibility under explicit code-based assumptions.


\subsection{C7 milestone status}

The C7 coordinate compactor is the conservative main candidate after the C6
relation-surface diagnostic.  It deliberately gives up the compactness of a full
projective block dictionary in exchange for removing the abundant within-block
weight-three projective relations.  This makes the current package ready as a
research prototype and manuscript milestone, while preserving a clear boundary:
it is not a production cryptosystem and does not certify concrete parameters.


\paragraph{Final v0.9-C7 status.}
The final validation pass moves the main compaction candidate from C4 projective
blocks to C7 coordinate relation-resistant compaction.  This choice sacrifices
the one-term-per-active-block advantage of C4, but removes the complete
projective weight-three relation surface discovered by C6.  The project is
therefore ready to pause as a research prototype and manuscript package, not as
a deployable cryptosystem.


\appendix
\section{Implementation-oriented pseudocode}
\label{app:pseudocode}

This appendix rewrites the construction in implementation-oriented form.  Field operations are modulo prime $q$.

\subsection{Noise and sparse sampling}

\begin{algobox}
\AlgTitle{$\mathsf{SampleNoise}(q,\eta)$}
\begin{enumerate}
    \item With probability $1-\eta$, return $0$.
    \item Otherwise return a uniformly random element of $\F_q^*$.
\end{enumerate}
\end{algobox}

\begin{algobox}
\AlgTitle{$\mathsf{SampleSparse}(n,k,q)$}
\begin{enumerate}
    \item Choose a uniformly random $k$-subset $S\subset[n]$.
    \item For each $i\in S$, choose $a_i\leftarrow\F_q^*$.
    \item Set all other coordinates to zero.
    \item Return sparse vector $a$.
\end{enumerate}
\end{algobox}

\subsection{Compact encryption}

\begin{algobox}
\AlgTitle{$\Regev\Enc_t(\mu)$}
\begin{enumerate}
    \item $a\leftarrow\mathsf{SampleSparse}(n,k,q)$.
    \item $e\leftarrow\mathsf{SampleNoise}(q,\eta)$.
    \item $b\leftarrow a\transpose t+\mu+e\pmod q$.
    \item Return $c=(a,b)$.
\end{enumerate}
\end{algobox}

\subsection{GSW-style scalar encryption}

\begin{algobox}
\AlgTitle{$\GSW\Enc_s(\alpha)$}
\begin{enumerate}
    \item For $j=1,\ldots,N$:
    \begin{enumerate}
        \item sample $a_j\leftarrow\mathsf{SampleSparse}(n,k,q)$;
        \item sample $e_j\leftarrow\mathsf{SampleNoise}(q,\eta)$;
        \item set row $z_j=(a_j,a_j\transpose s+e_j)$.
    \end{enumerate}
    \item Let $Z$ be the matrix with rows $z_j$.
    \item Return $X=Z+\alpha I_N$.
\end{enumerate}
\end{algobox}

\subsection{Code-LPN encryption}

\begin{algobox}
\AlgTitle{$\CLPN\Enc_r(x)$}
\begin{enumerate}
    \item Sample $A\leftarrow\F_q^{m_c\times n_c}$.
    \item Sample $e\leftarrow\chi_{q,\eta_c}^{m_c}$.
    \item Set $b=Ar+\mathsf{Code}(x)+e$.
    \item Return $(A,b)$.
\end{enumerate}
\end{algobox}

\begin{algobox}
\AlgTitle{$\CLPN\Dec_r(A,b)$}
\begin{enumerate}
    \item Compute $y=b-Ar$.
    \item Return $\mathsf{Decode}(y)$.
\end{enumerate}
\end{algobox}

\subsection{Key generation and evaluation}

\begin{algobox}
\AlgTitle{$\KeyGen$}
\begin{enumerate}
    \item Sample $t\leftarrow\F_q^n$, $s\leftarrow\F_q^n$, $r\leftarrow\F_q^{n_c}$.
    \item Set $\tilde{t}=(-t,1)$ and $\tilde{s}=(-s,1)$.
    \item For $i=1,\ldots,N$, set $\mathsf{EK}^{\rm exp}_i=\GSW\Enc_s(\tilde{t}_i)$.
    \item For $i=1,\ldots,N$, set $\mathsf{EK}^{\rm cmp}_i=\CLPN\Enc_r(\tilde{s}_i)$.
    \item Return $\sk=(t,s,r)$ and $\ek=(\mathsf{EK}^{\rm exp},\mathsf{EK}^{\rm cmp})$.
\end{enumerate}
\end{algobox}

\begin{algobox}
\AlgTitle{$\Expand(c)$}
\[
    C\leftarrow\sum_{i=1}^N c_i\mathsf{EK}^{\rm exp}_i.
\]
Return $C$.
\end{algobox}

\begin{algobox}
\AlgTitle{$\Compact(C)$}
\begin{enumerate}
    \item Let $\rho$ be the last row of $C$.
    \item Return
    \[
        \sum_{i=1}^N \rho_i\mathsf{EK}^{\rm cmp}_i.
    \]
\end{enumerate}
\end{algobox}

\section{Proof of the q-ary piling-up formula}
\label{app:piling-up}

We prove Lemma~\ref{lem:qary-piling}.  Let $E$ be a q-ary symmetric random variable with error probability $\eta$.  Let $\omega$ be a nontrivial additive character of $\F_q$.  Its Fourier transform is
\[
    \mathbb{E}[\omega(aE)]
    =1-\eta+\sum_{x\in\F_q^*}\frac{\eta}{q-1}\omega(ax)
\]
for $a\ne0$.  Since $\sum_{x\in\F_q}\omega(ax)=0$, we have $\sum_{x\in\F_q^*}\omega(ax)=-1$.  Therefore
\[
    \mathbb{E}[\omega(aE)]
    =1-\eta-\frac{\eta}{q-1}
    =1-\frac{q\eta}{q-1}.
\]
Multiplication by a nonzero coefficient $\alpha\in\F_q^*$ preserves the q-ary symmetric distribution.  For independent variables $E_1,\ldots,E_B$, the nontrivial Fourier coefficient of
\[
    S=\sum_{i=1}^B\alpha_iE_i
\]
is
\[
    \left(1-\frac{q\eta}{q-1}\right)^B.
\]
A q-ary symmetric variable with error probability $\eta_B$ has nontrivial Fourier coefficient
\[
    1-\frac{q\eta_B}{q-1}.
\]
Equating the coefficients gives
\[
    1-\frac{q\eta_B}{q-1}
    =\left(1-\frac{q\eta}{q-1}\right)^B,
\]
and hence
\[
    \eta_B=\frac{q-1}{q}\left(1-\left(1-\frac{q\eta}{q-1}\right)^B\right).
\]
This proves the formula.

\section{Reviewer and cryptanalysis checklist}
\label{app:checklist}

Before submitting a paper or releasing software, the following questions should be answered.

\subsection{Assumptions}

\begin{enumerate}
    \item Are the exact sparse-LPN parameters stated, including sample counts from all keys and ciphertexts?
    \item Does the proof use only standard sparse-LPN/q-ary LPN assumptions, or does it introduce an additional hidden assumption?
    \item Are low-noise regimes justified against known attacks?
    \item Does any optimization introduce ring, module, cyclic, quasi-cyclic, or ideal structure that changes the assumption?
\end{enumerate}

\subsection{Correctness}

\begin{enumerate}
    \item Are empirical failure rates consistent with the theoretical bounds?
    \item Is the final row support measured rather than only upper-bounded?
    \item Does the compactor decoder succeed at the observed effective q-ary error rate?
    \item Is repetition/plurality independent across replicas?
\end{enumerate}

\subsection{Security proof}

\begin{enumerate}
    \item Are $t$, $s$, and $r$ independent?
    \item Does the expansion key encrypt $\tilde{t}$ under $s$ only, avoiding circularity?
    \item Does the compaction key encrypt $\tilde{s}$ under $r$ only, avoiding circularity?
    \item Are all hybrids efficient and simulatable?
    \item Is the total exposed sample count within the assumed hardness regime?
\end{enumerate}

\subsection{Implementation}

\begin{enumerate}
    \item Are field operations constant-time where needed?
    \item Is noise sampling unbiased?
    \item Are sparse rows stored canonically to avoid duplicate-index bugs?
    \item Are decoding failures handled without leaking side information?
    \item Are parameter files machine-readable and bound to test vectors?
\end{enumerate}

\section{Attack-Screening Addendum}
\label{app:attack-screening}

This appendix records the first validation step after the initial construction: a conservative attack-screening layer for parameter rejection.  The estimates here are not security proofs and do not certify concrete SABLE parameters.  Their purpose is to expose weak settings before implementation or benchmarking.  In particular, any final parameter set must still be analyzed by specialists in sparse-LPN, q-ary LPN, information-set decoding, and large-sample attacks.

\subsection{Screened attack surfaces}

The validation repository models two public noisy-linear-equation surfaces.  First, the expansion key exposes approximately
\[
  m_{\exp}=(n+1)^2
\]
sparse-LPN-like rows related to the GSW expansion layer.  Second, the compaction key exposes
\[
  m_{\mathrm{cmp}}=(n+1)m_c
\]
q-ary LPN-like samples for the code/LPN compactor.  These are only proxy surfaces; the exact joint distribution of the GSW rows and encrypted secret-key coordinates must still be handled in the formal proof and in cryptanalysis.

The repository now implements three first-pass attack screens.

\paragraph{No-clean-subset screen.}
A random set of $n$ noisy equations is clean with probability roughly $(1-\eta)^n$.  The corresponding infinite-sample work proxy is
\[
  \log_2 W_{\mathrm{clean}}
  \approx -n\log_2(1-\eta)+\omega\log_2 n,
\]
where the last term represents dense linear algebra.  This screen is deliberately simple and is especially important when correctness pushes $\eta$ to be very small.

\paragraph{Prange/information-set screen.}
Viewing the public samples as a noisy random linear codeword, a Prange-style information-set decoder succeeds when it chooses an error-free information set.  With $m$ samples and expected error weight $t=\eta m$, the proxy work is
\[
  \log_2 W_{\mathrm{Prange}}
  \approx
  \log_2 {m\choose n} - \log_2 {m-t\choose n} + \omega\log_2 n.
\]
This is the most basic ISD screen, following the information-set viewpoint of Prange and later code-decoding attacks~\cite{Prange1962,Stern1989,BerlekampMcEliecevanTilborg1978}.

\paragraph{BKW-style bias screen.}
For q-ary symmetric noise, the estimator uses the character-bias proxy
\[
  \beta_q = \left|1-\frac{q\eta}{q-1}\right|.
\]
After approximately $2^a$ combined samples, the proxy bias is $\beta_q^{2^a}$ and a final distinguisher wants on the order of $\beta_q^{-2^{a+1}}$ samples.  The code minimizes a simple table-building and distinguishing proxy over the number of BKW levels.  This follows the role of BKW-family algorithms in LPN cryptanalysis~\cite{BlumKalaiWasserman2003,LevieilFouque2006,BogosEtAl2015}.  It is not a specialized sparse-q-ary-LPN estimator.

\subsection{Early finding: correctness/security tension}

The first estimator run rejects the earlier \texttt{prototype\_medium} preset.  The reason is not a software bug: the preset used very low noise for correctness.  For $q=65537$, $n=512$, $k=3$, and $\eta=2^{-20}$, the expansion-key surface has expected total errors below one.  The no-clean-subset and Prange-style proxies therefore fall far below a 128-bit target.  The compaction-key surface has the same problem.

The repository records the following conservative feasibility grid.  Here $\eta$ is set to the largest value allowed by the conservative requirement that the evaluated GSW error be at most $0.10$ for one additive term.  The column $n_{\min}$ is the smallest dimension making the no-clean-subset proxy reach 128 bits.  The final column is the sparse expansion-key entry count $(n+1)^2(k+1)$.

\begin{center}
\begin{tabular}{rrrrr}
\toprule
$k$ & depth & $\eta$ ceiling & $n_{\min}$ & expansion entries \\
\midrule
1 & 1 & $6.67\cdot 10^{-3}$ & $9{,}172$ & $1.68\cdot 10^8$ \\
2 & 1 & $2.50\cdot 10^{-3}$ & $23{,}388$ & $1.64\cdot 10^9$ \\
3 & 1 & $1.18\cdot 10^{-3}$ & $47{,}905$ & $9.18\cdot 10^9$ \\
1 & 2 & $3.92\cdot 10^{-4}$ & $135{,}802$ & $3.69\cdot 10^{10}$ \\
2 & 2 & $3.05\cdot 10^{-5}$ & $1{,}510{,}083$ & $6.84\cdot 10^{12}$ \\
3 & 2 & $4.58\cdot 10^{-6}$ & $8{,}895{,}213$ & $3.16\cdot 10^{14}$ \\
\bottomrule
\end{tabular}
\end{center}

The table suggests that the current construction is most plausible for depth-one arithmetic, i.e., degree-two computations, unless a new refresh, coding, or compaction technique improves the correctness/security tradeoff.  Depth two is not ruled out mathematically, but the naive expansion-key size becomes very large under this conservative screen.

\subsection{Design consequence}

The next research problem is therefore sharper than the original construction problem:

\begin{quote}
Can SABLE reduce its dependence on extremely low LPN noise while preserving low-depth homomorphic correctness?
\end{quote}

Promising directions are: using a stronger q-ary code rather than repetition in the compactor; reducing the public sample exposure through seeded or punctured evaluation keys; adding a lightweight code-based refresh for the GSW layer; deriving sharper circuit-aware error bounds instead of worst-case balanced-product bounds; and testing whether the sparse-LPN assumption used by Corrigan-Gibbs, Henzinger, Kalai, and Vaikuntanathan~\cite{CorriganGibbsHenzingerKalaiVaikuntanathan2024} admits concrete parameter regimes compatible with code/LPN compaction.

\section{C2 Block-Dictionary Compaction}
\label{app:c2-block-dictionary}

This appendix records the C2 refinement implemented in the validation repository.  The original compaction layer publishes one linearly homomorphic code-LPN encryption of each coordinate of the extended GSW secret $\tilde{s}\in\F_q^N$.  Compacting a final row $r\in\F_q^N$ then evaluates
\[
  \sum_{i=1}^N r_i \cdot \Enc(\tilde{s}_i),
\]
so the number of accumulated compaction-noise terms is the row support $B=\|r\|_0$.  Since q-ary symmetric noise piles up as
\[
  \eta_B = \frac{q-1}{q}\left(1-\left(1-\frac{q\eta_c}{q-1}\right)^B\right),
\]
large $B$ forces a very small compaction noise rate $\eta_c$.

The C2 block-dictionary compactor trades public-key size for fewer accumulated compaction-noise terms.  Partition $\tilde{s}$ into blocks
\[
  \tilde{s}=(\tilde{s}^{(1)},\ldots,\tilde{s}^{(L)}),
  \qquad \tilde{s}^{(j)}\in\F_q^{\ell_j},
\]
where $\ell_j\leq \ell$ and $L=\lceil N/\ell\rceil$.  For every block $j$ and every nonzero coefficient vector $u\in\F_q^{\ell_j}$, publish
\[
  \mathsf{EK}^{\mathsf{C2}}_{j,u}=\CLPN\Enc_r(\langle u,\tilde{s}^{(j)}\rangle).
\]
Given a final GSW row $r$, write it in matching blocks $r=(r^{(1)},\ldots,r^{(L)})$ and output
\[
  \mathsf{Compact}_{\mathsf{C2}}(r)
  =\sum_{j:r^{(j)}\neq 0}\mathsf{EK}^{\mathsf{C2}}_{j,r^{(j)}}.
\]
By linear homomorphism of the CLPN layer, the plaintext inside the compacted ciphertext is
\[
  \sum_{j:r^{(j)}\neq 0}\langle r^{(j)},\tilde{s}^{(j)}\rangle
  =r\cdot\tilde{s}.
\]
Thus C2 preserves the compaction semantics.

The number of compaction-noise terms becomes
\[
  t_{\mathsf{C2}}(r)=\#\{j:r^{(j)}\neq 0\}
  \leq \min\{\|r\|_0,\lceil N/\ell\rceil\}.
\]
For dense rows this improves the compaction-noise term by about a factor of $\ell$.  For sparse rows there may be no improvement.  The public-key cost is
\[
  \sum_{j=1}^L(q^{\ell_j}-1)
\]
CLPN ciphertexts, so this construction is only plausible for small $q$ and small $\ell$, or for CRT-laned designs over small prime fields.  This is consistent with the code-based secure aggregation literature, where CRT-style decompositions are used to reduce communication in LPN-based homomorphic aggregation \cite{BitzerEggerLiuWachterZeh2026}.

The security proof remains a hybrid argument from CLPN semantic security: the evaluator sees additional CLPN encryptions of linear forms of $\tilde{s}$, but these messages are hidden under the compaction secret.  The important change is quantitative rather than syntactic: C2 increases the number of public CLPN samples, so the attack surface must be re-estimated.  In the validation repository the C2 estimator therefore reports both the reduced compaction-noise count and the enlarged public sample surface.

The executable validation reports that, for the toy setting $q=7$, $N=13$, and $\ell=2$, the block dictionary contains 294 ciphertexts and correctly evaluates clean multiplication and $xy+z$ examples.  For the larger design-screening setting $q=7$, $N=513$, and $\ell=3$, the dictionary contains 58,482 CLPN ciphertexts.  At multiplicative depth two, the worst-case compaction terms drop from 513 to 171, but the public sample surface becomes very large.  Therefore C2 should be treated as a research refinement, not as a certified parameter set.

\section{C3 Seeded Block-Dictionary Storage}
\label{app:c3-seeded}

The C2 block-dictionary compactor reduces the number of accumulated CLPN noise
terms by publishing encryptions of all nonzero block inner products
$\langle u,\tilde{s}^{(j)}\rangle$.  Its direct implementation, however, stores a
full CLPN matrix $A$ and vector $b$ for every dictionary entry.  This appendix
records a storage refinement implemented in the validation repository.

\paragraph{Seeded CLPN ciphertexts.}
For each dictionary entry, sample a public seed $\sigma$ and define
$A_\sigma\in\F_q^{m_c\times n_c}$ by a deterministic expander.  Instead of
storing $(A_\sigma,b)$, the public key stores
\[
  (\sigma,b),
  \qquad
  b=A_\sigma r + x\mathbf{1}+e,
\]
where $x=\langle u,\tilde{s}^{(j)}\rangle$ is the dictionary plaintext and $r$
is the CLPN compaction secret.  Linear homomorphic addition concatenates seeded
matrix terms and adds the corresponding $b$ vectors.  To decrypt an aggregate
ciphertext, the holder of $r$ regenerates the required rows of each
$A_\sigma$ and plurality-decodes the residual.

\paragraph{Correctness.}
Let a compacted C3 ciphertext be the linear combination
\[
  (b,\{(\alpha_t,\sigma_t)\}_{t=1}^T),
  \qquad
  b=\sum_{t=1}^T\alpha_t b_t .
\]
Then
\[
  b-\sum_{t=1}^T\alpha_t A_{\sigma_t}r
  =\left(\sum_{t=1}^T\alpha_t x_t\right)\mathbf{1}
   +\sum_{t=1}^T\alpha_t e_t .
\]
Thus the plaintext and noise law are exactly those of dense C2.  Seeded storage
changes materialization cost but not the correctness invariant.

\paragraph{Security accounting.}
The public matrices are derivable from the seeds, so C3 exposes the same q-ary
LPN sample surface as dense C2: if the dictionary contains $M_{\mathrm{dict}}$ entries,
then the public compaction surface contains $M_{\mathrm{dict}}m_c$ LPN rows.  Therefore
seeding must not be claimed as a security improvement.  It is a key-size and
memory-traffic optimization only.  The validation repository adds a dedicated
q-ary/sparse-LPN surface estimator that reports the public sample count,
sample-to-dimension ratio, no-clean-subset and Prange-style screens, and a
q-ary block-BKW scan following the classical BKW line of work
\cite{BlumKalaiWasserman2003,LevieilFouque2006,BogosEtAl2015,EsserEtAl2017LPNDecoded}.

\paragraph{Storage comparison.}
Ignoring small serialization overheads, dense C2 stores approximately
\[
  M_{\mathrm{dict}}m_c(n_c+1)
\]
field elements.  Seeded C3 stores approximately
\[
  M_{\mathrm{dict}}m_c + M_{\mathrm{dict}}\cdot\frac{\lambda_{\mathrm{seed}}}{\log_2 q}
\]
field-element equivalents, where $\lambda_{\mathrm{seed}}$ is the public seed length in
bits.  The saving is therefore close to a factor of $n_c+1$ when $m_c$ is large
relative to the seed length.  The public attack surface remains unchanged.

\paragraph{Implemented status.}
The validation package includes seeded CLPN encryption, seeded block-dictionary
key generation, end-to-end clean multiplication tests, and a seeded C3 estimator.
The implementation is still a research prototype and uses SHA-256-derived rows
for reproducibility; a production design would require a specified, reviewed,
constant-time PRG or XOF interface and independent cryptanalysis.

\section{C2/C3 Public-Sample Surface Screen}
\label{app:c2-public-surface}

C2 and C3 reduce final compaction noise by replacing coordinate-wise compaction with a block dictionary.  The security cost is a much larger public q-ary LPN surface.  This appendix records the screening model used by the validation repository.  It is not a certified cryptanalytic estimator; it is a conservative design screen intended to expose surfaces that must be analyzed before any concrete security claim.

\paragraph{Dictionary row count.}
For block widths $\ell_j$, the C2 dictionary publishes
\[
  E = \sum_j (q^{\ell_j}-1)
\]
CLPN ciphertexts.  Each ciphertext contains $m_c$ noisy rows under the same compaction secret $r\in\mathbb F_q^{n_c}$, so the direct dictionary surface contains $E m_c$ public rows.  Seeded C3 changes only the representation of the public matrices: it stores a seed for each matrix and the corresponding $b$ vector.  Since the adversary can regenerate the public matrices from the seeds, the LPN sample count is unchanged.

\paragraph{Within-entry row differences.}
A CLPN dictionary entry has rows
\[
  b_i = \langle A_i,r\rangle + x + e_i \pmod q,
\]
where $x$ is the encrypted block inner product.  Subtracting two rows from the same entry cancels $x$:
\[
  b_i-b_j = \langle A_i-A_j,r\rangle + (e_i-e_j) \pmod q.
\]
Thus each dictionary entry yields up to $\binom{m_c}{2}$ message-eliminating q-ary LPN-like rows.  The effective noise rate is estimated by the q-ary piling-up formula with two terms,
\[
  \eta^{(2)} = \frac{q-1}{q}\left(1-\left(1-\frac{q\eta_c}{q-1}\right)^2\right).
\]
The validation code records the raw within-entry surface as
\[
  E\binom{m_c}{2}.
\]

\paragraph{Cross-entry dictionary differences.}
Inside one block, two dictionary entries for coefficient tuples $u$ and $v$ encrypt
\[
  \langle u,s_{\mathrm{block}}\rangle,\qquad \langle v,s_{\mathrm{block}}\rangle.
\]
A row difference between these entries gives
\[
  b_{u,i}-b_{v,j}
  = \langle A_{u,i}-A_{v,j},r\rangle
  + \langle u-v,s_{\mathrm{block}}\rangle
  + e_{u,i}-e_{v,j}.
\]
This is a noisy linear equation in the joint secret $(r,s_{\mathrm{block}})$ of dimension at most $n_c+\max_j\ell_j$.  The raw ordered surface is conservatively screened using
\[
  \sum_j \binom{q^{\ell_j}-1}{2} m_c^2
\]
rows, again with two-term q-ary piling-up noise.  These rows are not all independent, so the repository reports them as a red-flag screening surface rather than a proof of insecurity.

\paragraph{Implication for C3.}
Seeded C3 is still valuable because it avoids materializing $E m_c n_c$ public field elements for dense matrices and keeps only $E m_c$ field elements for the $b$ vectors plus seed metadata.  However, seeded storage does not reduce the within-entry or cross-entry surfaces above.  Therefore the next construction-level problem is not storage but relation control: a publishable C4 variant should avoid a fully closed block dictionary, restrict tuple sets, randomize dictionary availability, or replace the compactor with a code/LPN LHE layer whose public entries do not create so many plaintext-cancelling row differences.

\paragraph{Validation artifact.}
The repository implements this screen in the module
\begin{center}
\path{src/sable/c2_attack_surface.py}.
\end{center}
The generated toy and design reports in the documentation folder show that the current presets remain research-only.  This is expected and useful: the screen identifies the next mathematical bottleneck rather than certifying parameters.

\section{C4 Projective Sparse Additive-Basis Compaction}
\label{app:c4-projective}

This appendix records the C4 compaction refinement implemented in the validation repository.  The purpose of C4 is to reduce the public surface of the C2/C3 full block dictionary while preserving the small compaction-noise count that made block compaction attractive.

\subsection{Projective dictionary idea}

Let a final GSW last-row block be
\[
  \alpha = (\alpha_1,\ldots,\alpha_b) \in \F_q^b,
\]
and let the corresponding block of the extended GSW secret be
\[
  s_B = (s_1,\ldots,s_b) \in \F_q^b.
\]
A full C2 dictionary stores an encryption of
\[
  \langle \alpha,s_B\rangle
\]
for every nonzero \(\alpha\in\F_q^b\).  Hence one block requires
\[
  q^b-1
\]
public CLPN entries.

C4 observes that the linearly homomorphic compactor can scale ciphertexts.  Therefore it is unnecessary to publish separate entries for all scalar multiples of the same coefficient vector.  Let
\[
  \mathbb P^{b-1}(\F_q)
\]
denote the set of one-dimensional subspaces of \(\F_q^b\).  For each projective line, choose one representative \(u\) whose first nonzero coordinate is normalized to one.  The number of representatives is
\[
  \left|\mathbb P^{b-1}(\F_q)\right|
  = \frac{q^b-1}{q-1}.
\]
For every representative \(u\), the public C4 key stores
\[
  \Enc_{\mathrm{CLPN}}(\langle u,s_B\rangle).
\]

For any nonzero \(\alpha\in\F_q^b\), there exists a unique scalar \(\lambda\in\F_q^\ast\) and a unique normalized representative \(u\) such that
\[
  \alpha = \lambda u.
\]
The evaluator then computes
\[
  \lambda\cdot \Enc_{\mathrm{CLPN}}(\langle u,s_B\rangle),
\]
which is a CLPN encryption of
\[
  \lambda\langle u,s_B\rangle = \langle \alpha,s_B\rangle.
\]

\subsection{C4 correctness theorem}

\begin{theorem}[Projective C4 block correctness]
Let \(q\) be prime and let \(B\subseteq\{1,\ldots,N\}\) be a block of width \(b\).  Suppose the C4 public key contains CLPN encryptions of \(\langle u,s_B\rangle\) for every projective representative \(u\in\mathbb P^{b-1}(\F_q)\).  Then for every coefficient block \(\alpha\in\F_q^b\), C4 homomorphically produces a CLPN encryption of \(\langle \alpha,s_B\rangle\) using zero CLPN terms if \(\alpha=0\), and one CLPN term otherwise.
\end{theorem}

\begin{proof}
If \(\alpha=0\), the inner product is zero and the compactor contributes the zero CLPN ciphertext.  If \(\alpha\ne 0\), let \(i\) be the first index for which \(\alpha_i\ne 0\).  Set
\[
  \lambda = \alpha_i
  \quad\text{and}\quad
  u = \lambda^{-1}\alpha.
\]
Then the first nonzero coordinate of \(u\) is one, so \(u\) is exactly the projective representative stored by the public key.  By CLPN linear homomorphism,
\[
  \lambda\cdot \Enc_{\mathrm{CLPN}}(\langle u,s_B\rangle)
	ag{1}
\]
is a CLPN encryption of
\[
  \lambda\langle u,s_B\rangle
  = \langle \lambda u,s_B\rangle
  = \langle \alpha,s_B\rangle.
\]
This proves the claim.
\end{proof}

\subsection{Public-surface comparison}

For a block of width \(b\), the C4 projective dictionary has
\[
  M_{\mathrm{C4}}(b,q)=\frac{q^b-1}{q-1}
\]
entries, while C2/C3 full dictionaries have
\[
  M_{\mathrm{C2}}(b,q)=q^b-1.
\]
Thus
\[
  \frac{M_{\mathrm{C4}}}{M_{\mathrm{C2}}}=\frac{1}{q-1}.
\]
The compaction-noise term count remains one per active block, as in C2/C3.  Therefore C4 is not a storage-only transformation: it reduces both materialized key size and public CLPN sample count relative to full block dictionaries.

The cost is that C4 still publishes a structured family of related encryptions, one for every projective line in each block.  This does not by itself prove insecurity, but it is a public relation surface that must be analyzed by sparse-LPN and q-ary-LPN attack screens.

\subsection{Validation status}

The validation repository implements:
\begin{itemize}
  \item projective representatives and sparse decomposition;
  \item random additive-basis coverage screens;
  \item C4 key generation, compaction, and decryption;
  \item C4 public-surface estimates against v1, C2, and C3;
  \item end-to-end toy arithmetic-circuit tests.
\end{itemize}

The implementation should be interpreted as a research validation prototype.  The C4 construction improves the compaction-key public surface over C2/C3, but it does not remove the need for full LPN attack-cost estimation, large-sample analysis, and independent cryptanalysis.

\section{C5 Arithmetic Operation Suite and External Baseline Plan}
\label{app:c5-arithmetic}

The core construction is not limited to multiplication.  Once a compact input ciphertext is expanded into the GSW-style representation, the encrypted message space is the prime field $\F_q$.  If $C\tilde{s}=\mu\tilde{s}+e$, then linear public operations are immediate and nonlinear multiplication is matrix multiplication.

\paragraph{Native field operations.}
For expanded ciphertexts $C,C'$ and public $a,b\in\F_q$:
\begin{align*}
\Enc(0) &:= 0,\\
\Enc(b) &:= bI,\\
\Enc(\mu)+\Enc(\nu) &:= C+C',\\
\Enc(\mu)-\Enc(\nu) &:= C-C',\\
 a\Enc(\mu) &:= aC,\\
\Enc(\mu)\Enc(\nu) &:= CC'.
\end{align*}
The identity construction is exact because $(bI)\tilde{s}=b\tilde{s}$.  These operations support addition, subtraction, negation, public-scalar multiplication, public-constant addition, multiplication, squaring, public powers, and arbitrary low-degree public polynomials by composition.

\paragraph{Division and inversion.}
Encrypted division is not a native operation.  For a nonzero field element $x\in\F_q^*$, one may evaluate
\[
  x^{-1}=x^{q-2}
\]
by Fermat's little theorem.  This consumes multiplicative depth $O(\log q)$ by square-and-multiply and is undefined at $x=0$.  Therefore the prototype exposes nonzero inversion only as a validation tool for small $q$, not as a recommended primitive.

\paragraph{Boolean gates.}
Bits are encoded as $0,1\in\F_q$.  The validation code tests the polynomial gates
\begin{align*}
\mathsf{AND}(x,y)&=xy,\\
\mathsf{OR}(x,y)&=x+y-xy,\\
\mathsf{NOT}(x)&=1-x,\\
\mathsf{XOR}(x,y)&=x+y-2xy.
\end{align*}
Since the construction uses odd prime fields, XOR is not linear in this encoding.  This is a design tradeoff of the current q-ary sparse-LPN formulation.

\paragraph{Tested operations.}
Version C5 adds tests for constants, zero, one, addition, subtraction, negation, public-scalar multiplication, public-constant addition/subtraction, multiplication, oriented multiplication, square, powers, nonzero inverse, nonzero division, affine polynomials, and Boolean AND/OR/XOR/NOT/NAND/NOR/XNOR.

\paragraph{Baseline families.}
The external comparison plan separates exact arithmetic, approximate arithmetic, and Boolean/integer circuits.  OpenFHE documents support for BFV and BGV for integer arithmetic, CKKS for approximate real arithmetic, and FHEW/TFHE-style Boolean/arbitrary-function bootstrapping \cite{OpenFHEDocs2026}.  Microsoft SEAL supports BFV, BGV, and CKKS and is a natural exact/approximate arithmetic baseline \cite{MicrosoftSEAL2026}.  TFHE-rs targets Boolean, short-integer, and integer circuits and documents common integer arithmetic operations \cite{ZamaTFHERsDocs2026}.  The TFHE construction of Chillotti, Gama, Georgieva, and Izabachene remains the relevant Boolean-gate baseline \cite{CGGI2020}.

\paragraph{Interpretation of C5 measurements.}
The included timings are pure-Python prototype timings on toy parameters.  They validate operation semantics and expose relative costs inside the prototype.  They must not be interpreted as speed claims against optimized OpenFHE, SEAL, or TFHE-rs implementations.  The repository therefore includes an operation-support matrix and an optional external benchmark harness, but it does not fabricate wall-clock numbers for libraries that are not installed in the validation environment.

\section{C5 Baseline and Public-Surface Diagnostics}
\label{app:c5-baselines}

C5 separates measured prototype timing from baseline comparison.  The SABLE
repository is a pure-Python validation prototype, whereas mature libraries
such as OpenFHE, SEAL, Concrete, and TFHE-family implementations use
optimized C, C++, Rust, compiler passes, and architecture-specific kernels.
Therefore the current package reports operation-count and surface-area
proxies rather than claiming wall-clock superiority.

The fair baseline mapping is workload-dependent.  TFHE/FHEW-style schemes
are the closest comparison for Boolean gates and programmable function
evaluation \cite{DucasMicciancio2015,CGGI2020}.  BFV and BGV are the closest
comparison for exact modular or integer arithmetic \cite{FanVercauteren2012,BGV2014}.
CKKS is the closest comparison for approximate real or complex arithmetic,
but it is not the natural baseline for exact finite-field outputs
\cite{CKKS2017}.

The C5 public-surface screen also records that C4 projective compaction
removes two-term projective duplicates but necessarily introduces many
three-term linear relations once the block width is at least two.  For public
representatives \(u,v,w\) of projective lines, many choices satisfy
\[
    au+bv-cw=0
\]
for non-zero field elements \(a,b,c\).  This relation surface does not by
itself break the CLPN layer, but it does give an adversary many public noisy
linear combinations to examine.  Consequently, the C5 conclusion is that C4
is smaller than C2/C3, but a security-grade claim requires a dedicated
large-sample q-ary-LPN analysis of the projective public surface.

\section{C6 Projective-Relation Surface Estimator}
\label{app:c6-relation-surface}

C6 adds a dedicated attack-surface screen for the C4 projective compactor.
The C4 construction publishes one representative of each one-dimensional
subspace in a block.  For a block of width $b$ over $\F_q$, the number of
public representatives is
\[
    M_b = \frac{q^b-1}{q-1}.
\]
This is a factor of roughly $q$ smaller than the non-zero C2 dictionary, but
it is still a highly structured public set.

\paragraph{Projective-line relations.}
The important C6 correction is that a full projective block of width
$b\geq 2$ has many weight-three linear dependencies.  The number of
projective lines in $\mathrm{PG}(b-1,q)$ is the Gaussian binomial coefficient
\[
    {b \brack 2}_q
    =
    \frac{(q^b-1)(q^{b-1}-1)}{(q^2-1)(q-1)}.
\]
Every projective line has $q+1$ points, and every triple of distinct points
on that line is linearly dependent.  Thus the raw number of weight-three
projective-line relations in one block is
\[
    R_{3,b}
    =
    {b \brack 2}_q {q+1 \choose 3}.
\]
For $b=1$ this number is zero; for every $b\geq 2$ it is positive.  This
replaces the earlier rough C5 heuristic that treated the generic minimum
relation weight as $b+1$.  The generic statement is true for arbitrary
vectors in general position, but the complete projective set is not in
general position: it contains complete projective lines.

\paragraph{Known-zero CLPN combinations.}
Let the C4 key contain CLPN ciphertexts of
$\langle u_j,s_B\rangle$ for public block masks $u_j$.  If
\[
    a u_i + b u_j + c u_k = 0
\]
for non-zero $a,b,c\in\F_q$, then the same linear combination of the three
CLPN ciphertexts encrypts the known message zero.  The resulting row-level
samples have q-ary symmetric noise rate
\[
    \eta_3
    =
    \frac{q-1}{q}
    \left(1-
    \left(1-\frac{q\eta_c}{q-1}\right)^3
    \right),
\]
by the q-ary piling-up formula.  These relation-derived samples must be
screened separately from ordinary row differences inside a single CLPN
ciphertext, which have effective noise $\eta_2$.

\paragraph{Raw count versus rank-capped count.}
The raw number $R_{3,b}$ should not be interpreted as an independence count.
The space of all linear relations among the $M_b$ public projective
representatives has dimension
\[
    K_b = M_b-b.
\]
Therefore the C6 estimator reports both
\[
    m_c \sum_B R_{3,b_B}
\]
as a raw low-weight relation surface and
\[
    m_c \sum_B K_{b_B}
\]
as a rank-capped independence proxy.  Both quantities matter: the raw count
measures the adversary's search surface, while the rank-capped count prevents
the screen from pretending that all low-weight relations are independent.

\paragraph{C6 diagnostic verdict.}
The C6 validation package screens four surfaces:
\begin{enumerate}
    \item sparse-LPN expansion-key rows;
    \item C4 CLPN row-difference samples;
    \item C4 weight-three relation samples with rank-capped counting;
    \item C4 weight-three relation samples with raw upper-bound counting.
\end{enumerate}
The current C6 result rejects full projective C4 as the main security
candidate for block width at least two.  This is not a correctness failure;
it is a public-surface warning.  The algebraic compactor works, but abundant
weight-three public relations give an attacker known-zero q-ary-LPN-style
samples at low accumulated noise.  This interacts directly with the known
importance of BKW-style algorithms and information-set-decoding attacks for
LPN and code-based cryptography \cite{BlumKalaiWasserman2003,BogosEtAl2015,Prange1962,Stern1989,MayMeurerThomae2011,BeckerJouxMayMeurer2012}.

\paragraph{Implication for the next construction.}
C7 should replace the full projective dictionary with a relation-resistant
compaction layer.  Possible directions are incomplete randomized additive
bases with screened sparse-kernel weight, salted block masks that prevent
public cancellation to known-zero messages, or a decoding-based compactor
whose public generator has no abundant low-weight kernel.  The target
property is no longer only coverage; the public basis must also have high
minimum sparse-kernel weight while maintaining compaction noise low enough for
correctness.

\section{C7 Relation-Resistant Compaction and Final Validation Status}
\label{app:c7-relation-resistant}

C7 is the final construction pass in this validation package.  It responds to
C6's negative finding: complete projective C4 blocks are correct, but they
publish many weight-three public relations.  C7 therefore changes the main
candidate compactor from a full projective block dictionary to a
relation-resistant coordinate compactor, while keeping screened random masks as
an experimental optimization.

\subsection{C7 coordinate compactor}

Partition the final GSW last-row coefficient vector into blocks
$\alpha_B\in\F_q^b$.  In the C7 coordinate mode, the public basis for a block
is simply
\[
    U_B = \{e_1,\ldots,e_b\},
\]
where $e_i$ are the standard coordinate vectors.  The public key stores
\[
    \CLPN.\Enc_r(\langle e_i,s_B\rangle)
    =
    \CLPN.\Enc_r(s_{B,i})
\]
for each coordinate in the block.  Evaluation decomposes
\[
    \alpha_B=\sum_{i=1}^b \alpha_i e_i
\]
and linearly combines the corresponding CLPN ciphertexts.

This mode is intentionally conservative.  It does not reduce the worst-case
number of compaction terms below the legacy coordinate compactor, but it also
does not create the complete projective-line relation surface from C4.  In
particular, inside a block the matrix whose rows are $e_1,\ldots,e_b$ has full
row rank and no non-zero kernel relation.

\begin{lemma}[C7 coordinate coverage]
For every $\alpha_B\in\F_q^b$, the C7 coordinate basis gives a decomposition
of Hamming weight exactly $\wt(\alpha_B)$ and at most $b$.
\end{lemma}

\begin{proof}
The coordinate identity
$\alpha_B=\sum_i \alpha_i e_i$ is unique.  Zero coefficients contribute no
term, so the number of non-zero terms is $\wt(\alpha_B)\leq b$.
\end{proof}

\begin{lemma}[No within-block kernel relation]
The C7 coordinate basis has no non-zero linear relation among its public masks.
\end{lemma}

\begin{proof}
If $\sum_i z_i e_i=0$, then the $i$th coordinate of the sum is $z_i$.  Hence
all $z_i=0$.
\end{proof}

\subsection{Experimental screened-random mode}

C7 also includes an optional screened-random mode.  Each block begins with the
coordinate basis so that correctness is guaranteed.  The generator then tries
to add dense random masks.  A candidate mask is rejected if it is projectively
duplicate or if adding it creates a local kernel relation up to the configured
screen weight.  During compaction, the evaluator first searches for a sparse
representation using the extra masks and falls back to the coordinate
representation when no sparse representation is found.

This mode is not the main security claim.  It is an optimization target for
future work because every added mask increases the public LPN sample surface.
The final package therefore marks C7 coordinate as the main research default
and screened-random C7 as experimental.

\subsection{Correctness theorem}

Let $C_\star$ be the evaluated GSW ciphertext and let $r_\star$ be its last
row.  Let $B_1,\ldots,B_L$ be the coefficient blocks touched by $r_\star$.
C7 coordinate compaction produces
\[
    \sum_{B}\sum_{i\in B} r_{\star,i}\,\CLPN.\Enc_r(s_i),
\]
which is a CLPN encryption of
\[
    \sum_i r_{\star,i}s_i = r_\star\cdot \widetilde{s}.
\]
As in the main construction, if
\[
    r_\star\cdot\widetilde{s}=f(\mu_1,\ldots,\mu_m)+e_\star,
\]
then the compacted ciphertext decrypts to the correct value whenever
$e_\star=0$ and the CLPN decoder succeeds.

If $T$ CLPN ciphertexts are added during C7 compaction, the q-ary symmetric
error rate entering the repetition decoder is bounded by
\[
    \eta_T
    =
    \frac{q-1}{q}
    \left(
    1-
    \left(1-\frac{q\eta_c}{q-1}\right)^T
    \right).
\]
For coordinate C7, $T\leq \wt(r_\star)\leq N$.  This is worse than one-term
projective C4 in pure noise terms, but it removes the complete projective
weight-three relation surface.

\subsection{Validation results in the repository}

The v0.9-C7 package contains the following additional validation artifacts:
\begin{itemize}
    \item \texttt{src/sable/c7\_relation\_resistant.py}: C7 coordinate and screened-random compaction;
    \item \texttt{experiments/run\_c7\_arithmetic\_suite.py}: all arithmetic and Boolean-operation tests under C7;
    \item \texttt{experiments/run\_c7\_compare\_baselines.py}: operation-count comparison against TFHE/FHEW, BFV/BGV, and CKKS-style families;
    \item \texttt{experiments/run\_c7\_readiness.py}: final readiness report;
    \item \texttt{docs/c7\_readiness\_output.json}: machine-readable status summary.
\end{itemize}

The current C7 validation run reports that all tested arithmetic operations
succeed in the toy-clean configuration: addition, subtraction, negation,
public scalar multiplication, public constants, multiplication, squaring,
affine combinations, dot products, polynomial evaluation, balanced products,
quadratic forms, and the Boolean gates NOT, AND, OR, XOR, NAND, NOR, XNOR, and
implication encoded over $\F_q$.  The test suite reports \texttt{95 passed}.

\subsection{Final project status}

At the end of C7, the project is ready to pause as a research prototype and
manuscript package.  The main candidate is no longer C4 projective compaction;
it is C7 coordinate relation-resistant compaction.  The package is not ready
for deployment or standardization, and it does not claim certified concrete
security.  The remaining scientific work is independent cryptanalysis,
stronger q-ary decoding beyond repetition coding, and measured benchmarks
against optimized HE implementations on matched circuits.


\input{references_manual.tex}

\end{document}
